% =====================================================================
%  BAB IV HASIL DAN PEMBAHASAN  (khusus sidang akhir)
% =====================================================================
\chapter{HASIL DAN PEMBAHASAN}
\label{bab:hasil}

\petunjuk{Inti laporan: temuan dan interpretasinya (Pedoman TA 7.4). Revisi
yang mengubah substansi bab ini tergolong revisi sedang atau mayor.
\tahap{\TidakPerlu}{\WAJIB}}

\section{Hasil Penelitian}
\petunjuk{Sajikan hasil sesuai urutan tahapan metodologi dalam tabel, grafik,
atau diagram yang bernomor, berjudul, dan dirujuk di narasi.}

\isian{Tulis uraian dalam bentuk paragraf.}

\section{Analisis dan Pembahasan}
\petunjuk{Jelaskan makna dan pola data, bandingkan dengan penelitian terdahulu
(Bab 2.2), identifikasi faktor yang memengaruhi hasil, dan jawab setiap
rumusan masalah. Jangan sekadar mengulang isi tabel.}

\isian{Tulis uraian dalam bentuk paragraf.}

\section{Keterbatasan Penelitian}
\petunjuk{Akui secara jujur keterbatasan metode, dataset, lingkungan
eksperimen, dan cakupan analisis.}

\isian{Tulis uraian dalam bentuk paragraf.}
